% s-132-p132-m004; original PDF 132
\Aphorism{s-132-p132-m004}{186}
Man kommt sich als Liebender zu gut vor für diese Welt, man sieht sich \Foreign{sub specie aeterni,} man trägt einen Ernst zur Schau, der dem Wissenden von gestörter Verdauung (und nicht von „metaphysischem Bewusstsein“) erzählt, man frisirt sich \Foreign{malcontent,} — kurz man macht vor der Nachwelt, die man in sich beherbergt, eine komische Figur. Muss das sein? kann man nicht lieben und dabei obenauf bleiben, kann man nicht eine altfranzösische intelligente Liebesgeschichte erleben, deren episodische Ausdehnung von vornherein zu übersehen ist, etwas Nervenberuhigendes, wobei der gebildete Europäer sich selbst noch wiedererkennt? Aber die \Foreign{mater saeva cupidinum} hat Freude an unserer Instinctbenebelung und lässt den würdigsten Herrn, Doctor aller vier Facultäten, nach Gretchens Strumpfbande schmachten: „der grosse Hans, ach wie so klein \Dots{3}“
\par

